
\section{A Taxonomy of Self-Improvement Mechanisms}
\label{sec:A_Taxonomy_of_Self_Improvement_Mechanisms}

Building on the formal definition of the foundation-model-based agent $\mathcal{A}_{t}=(\theta_{t}, \Sigma_{t})$ introduced in Section \ref{subsec:Foundation_Model_Based_Agents}, where $\Sigma_{t} := (p_{t}, m_{t}, \mathcal{T}_{t}, g_{t})$, we now establish a unified taxonomy for self-improvement. 

Historically, as discussed in Section \ref{subsec:Self_Improvement_in_FM_based_Agents}, self-improvement was often viewed strictly through the lens of parameter optimization (e.g., connectionism) or code rewriting (e.g., Gödel Machines). However, the emergence of Large Language Models (LLMs) has bifurcated the improvement landscape into two distinct but complementary pathways: one targeting the \textit{internal cognitive substrate} (the model weights) and the other targeting the \textit{external operational framework} (the scaffolding).

% We formally define the self-improvement operator $\IMPROVE$ as a mapping txu2023wizardlmhat transitions the agent from state $t$ to $t+1$ based on a learning signal $\mathcal{S}_t$:
More formally, the self-improvement process is defined as a self-improvement operator, $\IMPROVE(\cdot;\cdot)$, which maps an agent's current state and a learning signal to a new, improved state. This operator acts on one or more components of the agent. The general formula is:
\begin{equation}
    \mathcal{A}_{t+1} = \IMPROVE(\mathcal{A}_{t}; \mathcal{S}_{t})
\end{equation}
Depending on which component of the tuple $(\theta_{t}, \Sigma_{t})$ is the target of the operator $\IMPROVE$, we classify self-improvement mechanisms into two primary paradigms: \textit{Foundation Model Improvement} (Parametric) and \textit{Scaffolding Improvement} (Non-Parametric).

This chapter will elaborate on these two core paradigms, which serve as a roadmap for the detailed technical discussions in the subsequent sections.

\subsection{Foundation Model Improvement (The Parametric Slow Loop)}
\label{subsec:Foundation_Model_Improvement_(The_Parametric_Slow_Loop)}

This paradigm represents the ``Internalization'' of capabilities. It targets the foundation model parameters $\theta_{FM}$, treating the agent's neural substrate as a fluid medium that can be reshaped by experience. 

In this mode, the agent directly updates the weights of its neural network through methods like fine-tuning, based on experience and feedback, thereby solidifying new knowledge or skills into its internal core. This can be expressed as:
\begin{equation}
    \theta_{t+1} = \IMPROVE_{\theta}(\theta_{t}; \mathcal{S}_{t}), \quad \Sigma_{t+1} = \Sigma_{t}
\end{equation}
Here, the scaffolding remains constant (or is updated separately), while the weights $\theta$ are updated via gradient-based optimization methods such as Supervised Fine-Tuning (SFT) or Reinforcement Learning (RL). This form of improvement corresponds to the \textbf{Slow Loop} of learning: it is computationally expensive and stable, intended to permanently encode knowledge, reasoning patterns, and alignment preferences into the agent's long-term memory (the weights).

As detailed in Section \ref{sec:Foundation_Model_Improvement}, we categorize this paradigm based on the source of the learning signal $\mathcal{S}_t$:
\begin{itemize}
    \item \textbf{Self-Generated Data ($\mathcal{S}_t \approx \mathcal{D}_{t}$):} The agent acts as a \textit{teacher} by synthesizing explicit training instances $\mathcal{D}_{t}$, such as high-quality demonstrations or augmented datasets, to expand its own knowledge boundaries via supervised learning ~\citep{wang2022self, lu2023self, zhao2024selfguidebettertaskspecificinstruction, shi2025taskcraft}.
    \item \textbf{Self-Generated Supervision ($\mathcal{S}_t \approx r_{t}$):} The agent acts as a \textit{critic} by constructing internal supervisory signals $r_{t}$\ in the form of scalar rewards, preference pairs, or verbal critiques to guide policy optimization and alignment \citep{gong2025strive, huang2025beyond, bai2022constitutional, bensal2025reflect, zhang2025adaptive}.
    \item \textbf{Self-Generated Experience ($\mathcal{S}_t \approx \tau_{t}$):} The agent acts as an \textit{explorer} by learning from environment-grounded outcomes $\tau_{t}$, where improvement is driven by interaction trajectories and feedback-rich traces encountered during task execution ~\citep{wei2025codearc, dong2025tool, da2025agent, sun2025seagent}.
\end{itemize}

\subsection{Scaffolding Improvement (The Non-Parametric Fast Loop)}
\label{subsec:Scaffolding_Improvement_(The_Non_Parametric_Fast_Loop)}

This paradigm represents the ``Externalization'' of capabilities. It targets the operational scaffolding $\Sigma$, enhancing the agent by optimizing the context, tools, and memory structures that surround the frozen foundation model. This can be expressed as:
\begin{equation}
    \Sigma_{t+1} = \IMPROVE_{\Sigma}(\Sigma_{t}; \mathcal{S}_{t}), \quad \theta_{t+1} = \theta_{t}
\end{equation}
This corresponds to the \textbf{Fast Loop} of learning: it is computationally efficient, highly adaptive, and reversible. It allows the agent to adapt to new tasks on-the-fly without the risk of catastrophic forgetting associated with parameter updates.

As detailed in Section \ref{sec:Scaffolding_Improvement}, we decompose this paradigm based on the specific component of $\Sigma_{t} := (p_{t}, m_{t}, \mathcal{T}_{t}, g_{t})$ being modified:

\begin{itemize}
    \item \textbf{Prompt Optimization ($p_{t} \rightarrow p_{t+1}$):} Improving the \textit{interface logic.} This paradigm targets the interaction level, where the agent refines its instructions and in-context demonstrations to elicit better reasoning from $\theta$. This process focuses on optimizing the conditioning context without altering underlying parameters, which is formally expressed as $p_{t+1}=\IMPROVE_{p}(p_t; \mathcal{S}_t)$ ~\citep{yuksekgonul2024textgradautomaticdifferentiationtext, guo2025evopromptconnectingllmsevolutionary, fernando2024promptbreeder, zhou2022large}.

    \item \textbf{Memory Evolution ($m_{t} \rightarrow m_{t+1}$):} Improving the \textit{experience management.} This mechanism enhances how the agent retains and structures its history, moving from passive logging to active consolidation. By reasoning over past interactions to ensure long-term coherence, the agent updates its memory state as $m_{t+1}=\IMPROVE_{m}(m_t; \mathcal{S}_t)$ ~\citep{chhikara2025mem0buildingproductionreadyai, wang2024agentworkflowmemory, zhang2025memgenweavinggenerativelatent, zhong2024memorybank}.

    \item \textbf{Tool Governance ($\mathcal{T}_{t} \rightarrow \mathcal{T}_{t+1}$):} Improving the \textit{action space.} This paradigm focuses on the agent's ability to transcend inherent limitations by creating, filtering, or refining the external APIs it invokes. By ``learning to use or build tools,'' the agent expands its functional scope, formalized as $\mathcal{T}_{t+1}=\IMPROVE_{\mathcal{T}}(\mathcal{T}_t; \mathcal{S}_t)$ ~\citep{haque2025advancedtoollearningselection, wang2025toolgenunifiedtoolretrieval, dong2025tool, zhang2025agentorchestraorchestratinghierarchicalmultiagent}.

    \item \textbf{Full Scaffolding Update ($\Sigma_{t} \rightarrow \Sigma_{t+1}$):} Improving the \textit{system architecture.} A holistic reconfiguration where the agent modifies its own core operational logic or source code. This involves a recursive self-improvement process that integrates prompts, memory, and tools into a new architectural state, represented as $\Sigma_{t+1}=\IMPROVE_{\Sigma}(\Sigma_t; \mathcal{S}_t)$ ~\citep{hu2025automateddesignagenticsystems, zhang2025darwingodelmachineopenended, novikov2025alphaevolvecodingagentscientific, xia2025livesweagentsoftwareengineeringagents}.
\end{itemize}


\subsection{The Unified Cycle}
\label{subsec:The_Unified_Cycle}

While we treat these paradigms separately for analytical clarity, it is conceptually useful to view them as two complementary loops that could form a unified cycle. In this idealized workflow, an agent first performs \textit{Scaffolding Improvement} to adapt quickly to a target task, e.g., refining prompts, updating memory, or adjusting tool usage, thereby exploring solutions under a fixed foundation model (the Fast Loop). The resulting interaction traces, verified solutions, and feedback signals can then be curated into self-generated datasets or supervision, which in turn drive parameter updates via fine-tuning or reinforcement learning (the Slow Loop), internalizing the acquired capability.

Crucially, this two-loop coupling separates \emph{transient adaptation} from \emph{permanent acquisition}. Scaffolding updates provide reversible, low-cost plasticity that enables rapid experimentation and safe iteration, while parametric updates provide stable, amortized learning that transfers across future tasks and contexts. This unified perspective offers a principled lens for understanding how modern agents evolve: they \emph{externalize} capabilities to solve problems efficiently in the moment, and then \emph{internalize} the verified improvements to expand their long-term competence.
